% s00-intro.tex —— 印刷页 1（引言）
% 引言在原书中无节号；此处用不编号节，避免占用第 1 节号而使 §1 几何类域论
% 及其后各节的定理编号整体后移一位。
\section*{引言}
\addcontentsline{toc}{section}{引言}
\srcpage{1}

1967 年，R. Langlands 提出了 E. Artin 与 J. Tate 的阿贝尔类域论的一个宏大扩展。更确切地说，他猜想：整体域 $F$ 上约化群 $G$ 的自守表示，与 $F$ 的取值于 $G$ 的对偶代数群 ${}^{L}G$ 的 Galois 表示之间，存在自然的对应。${}^{L}G$ 的单位分支 $\widehat{G}$ 是复约化群，其根是 $G$ 的余根，反之亦然。

若 $G$ 是 $F$ 上的线性群 $\mathrm{GL}(n)$，则 $\widehat{G}$ 正是 $\mathrm{GL}(n,\C)$。当 $F$ 是函数域，即 $\F_p(t)$ 的有限扩张时，整体 Langlands 对应已由 V. Drinfeld [Dr 1]（$n=2$）与 L. Lafforgue [La]（$n$ 任意）证明。数域上的整体 Langlands 对应仍是数学中尚未解决的大问题之一。

Drinfeld–Langlands 对应，又称几何 Langlands 对应，是 Langlands 对应在以下情形的猜想类比：$k(t)$ 的有限扩张 $F$ 上的分裂约化群 $G$，其中 $k$ 是任意域。若 $X$ 是 $k$ 上光滑拟射影的代数曲线，函数域为 $F$，则该对应使 $X$ 上 $G$-丛的模空间与 $X$ 上 $\widehat{G}$-局部系统的模空间互为对偶。

当 $G=\mathrm{GL}(1)$ 时，Drinfeld–Langlands 对应正是 M. Rosenlicht 与 S. Lang 的几何类域论，由 J.-P. Serre 在 [Se] 中阐述。$G=\mathrm{GL}(2)$ 的情形由 V. Drinfeld 在 [Dr 2] 与 [Dr 3] 两文中处理，这两篇文章正是该理论的源头。

许多工作研究了 Drinfeld–Langlands 对应的各个方面，特别是 A. Beilinson 与 V. Drinfeld（[B-D]）、A. Braverman 与 D. Gaitsgory（[B-G]）、以及 S. Lysenko（[Ly 1] 与 [Ly 2]）的工作。在本文中，我只讨论 E. Frenkel、D. Gaitsgory 与 K. Vilonen 的近期工作，即 $X$ 射影（处处非分歧对应）且 $G=\mathrm{GL}(n)$ 的情形。这些工作推广了 V. Drinfeld 在 [Dr 2] 中的结果。

我感谢 S. Lysenko 在本文准备过程中给予的帮助。
